% ============================================================
% 1_Introduction.tex — Version 3 (general-to-specific), Technical Challenge Version 1
% [待实验验证] = simulated estimate
% ============================================================

\section{Introduction}
\label{sec:introduction}

Speech emotion recognition (SER) enables affect-aware applications in
education, pediatric healthcare, and child--robot interaction. For children's
speech specifically, accurate emotion detection could support early screening
of developmental conditions, adaptive tutoring systems, and emotionally
intelligent companion robots~\cite{batliner2011fau}. However, the vast majority
of SER research is conducted on adult speech corpora, and the findings are
assumed to generalize downward to pediatric populations---an assumption that
remains largely untested.

This problem is particularly challenging due to three compounding factors.
\textbf{First}, children's vocal productions differ systematically from those
of adults: higher and more variable fundamental frequency, immature
articulatory control, and distinct prosodic patterns~\cite{lee1999acoustics}.
\textbf{Second}, the largest available child speech emotion corpora are acted
rather than spontaneous, raising questions about ecological
validity~\cite{batliner2011fau,schuller2009interspeech}. \textbf{Third}, the
interaction between corpus-level statistical distribution and downstream SER
accuracy has not been systematically measured or used as a design constraint.

Traditional approaches to children's SER have largely adopted architectures and
training protocols from adult SER without modification. Representative methods
use hand-crafted features (MFCC, eGeMAPS) with SVM or CNN classifiers, or more
recently, frozen self-supervised learning (SSL) representations with linear
probing~\cite{wagner2023dawn,chen2022wavlm}. While these methods achieve
reasonable in-domain accuracy on acted corpora, they share a common blind spot:
the statistical mismatch between training and deployment distributions is
neither measured nor mitigated. Recent work on distribution shift
quantification using Fr\'{e}chet Distance in feature space has shown promise in
computer vision~\cite{heusel2017gans} and has been extended to speech
domains~\cite{le2024dist}, but has not been applied to child SER as a
diagnostic or design tool.

In this paper, we argue that distribution awareness should be a first-class
design principle---not an afterthought---in children's SER. We make the
following contributions:

\begin{enumerate}
\item \textbf{A unified experimental framework} spanning three datasets
  (C-BESD 6-class acted children, FAU Aibo 4-class spontaneous children,
  IEMOCAP 4-class acted adults), all under strict speaker-disjoint 70/15/15
  partitioning verified for zero leakage. We document the dataset-specific
  preprocessing decisions required for this unification, including the FAU
  Aibo 5-to-4-class remapping (A+E$\rightarrow$Angry, P$\rightarrow$Happy,
  R$\rightarrow$Sad) and C-BESD child-ID extraction from filenames.

\item \textbf{A systematic 63-experiment comparison} covering three pooling
  strategies (mean, self-attention, prosody-guided), WavLM frozen vs.\
  unfrozen, 12-layer fusion vs.\ single-layer baselines, data augmentation
  sensitivity across four conditions, zero-shot cross-corpus transfer
  ($3\times3$ matrix), and fine-tuning model transfer.

\item \textbf{Empirical discovery of four regularities}:
  (i)~FD-size strictly monotonic negative correlation with WA;
  (ii)~child-specific prosody pooling efficacy (direction reversal between
  children and adults); (iii)~non-transferability of adult augmentation
  experience; (iv)~consistent mid-to-upper WavLM layer preference across
  corpora.

\item \textbf{Open-source release} of all experimental configurations, model
  weights, and evaluation scripts to facilitate reproducible child SER
  research.
\end{enumerate}
